% AI-written, AI cross-reviewed; no human sign-off; not a claim that RH is proved
% RH Notes blueprint source excerpts for Z4.2.04
% Excerpts in the notation of the parent blueprint. The local edition corrects the attribution of the round-16 and round-17 reviews. The labelled editions and source ranges appear below.
% See the unit YAML for the current proof status. Earlier-round passages retain their original wording.
\section*{Z4.2.04. Theta kernels for adjacent pairs}

% Parent source SHA-256: 1af64922b98ed2310744411b1f18b25cc426c97220646321471836e1a3a409d6
% blueprint/ch_P.tex lines 13181 through 13187; r6:Z4.2.04
\subsection{Z4.2.04. The two frames and the finite theta sum (v8)}\label{r6:Z4.2.04}
\emph{The v10 continuation is on page~\pageref{r8:Z4.2.04}.}
Fix \(j<m\) and valid \(L,L'\) at \(j,j+1\). Reindexing identifies \(N_L\) with the adelic lattice at a unitary point. The unitary action transports the self-dual standard lattice to a self-dual lattice, and reindexing preserves its dual; the rank at the site is even. Apply the same argument at \(j+1\). For the span calculation, choose a local basis for the tuple. Its columns span the prescribed local lattice. Mapping those columns by the frame gives the columns in the adelic lattice after taking the first factor. At a successor one first transports along the rank identification of that chart. Hence the span is the \emph{image} of the local lattice under the frame at each of the two ends. Both frame determinants are units (\nolinkurl{even_and_isSDLat_kerN}, \path{Development/AbelianVarietyTheory/Proofs/SatakeParameters/KernelLatticeSelfDualEvenRankAtSite.lean:13--29}; \nolinkurl{locSp_kerN_zero}, \path{Development/AbelianVarietyTheory/Proofs/SatakeParameters/FrameLocalLatticeSpanKernelLatticeZero.lean:19--35}; \nolinkurl{locSp_kerN_succ}, \path{Development/AbelianVarietyTheory/Proofs/SatakeParameters/FrameLocalLatticeSpanKernelLatticeSucc.lean:18--52}; \nolinkurl{GlobalTheory.pairOK_allSites}, \path{Development/AbelianVarietyTheory/Proofs/SatakeParameters/KernelSiteLatticeTuplePairsAdmissible.lean:25--47}).

Pull the finite datum back to the adelic charts. There its value depends only on the specified congruence class; its values are integral away from the specified bad primes. The admitted lattices give an integral, finitely generated period lattice. Passing to the positive reader level gives finite index, so the kernel is a finite sum \(\sum_i r_i U_{t_i}e_{\rm nat}\) with algebraic \(r_i\) and with a positive multiple of each \(t_i\) in the period lattice. Finite sums preserve differentiability and algebraicity of Taylor coefficients separately. The scalar and the nonzero algebraic character factor are supplied by the determinant-character calculation on this same pair (\nolinkurl{xiM_datLoc}, \path{Development/AbelianVarietyTheory/Proofs/SatakeParameters/MovedPairDatumDependsOnlyOnLevel.lean:13--30}; \nolinkurl{kerTh_entire_algCoef}, \path{Development/AbelianVarietyTheory/Proofs/SatakeParameters/KerThEntireAlgCoef.lean:22--40}; \nolinkurl{latticeTuplePairScalar_mem_algebraicClosure_alt}, \path{Development/AbelianVarietyTheory/Proofs/SatakeParameters/LatticeTuplePairScalarMemAlgebraicClosure.lean:13--26}; \nolinkurl{pairOK_allSites}, \path{Development/AbelianVarietyTheory/Proofs/SatakeParameters/KernelSiteLatticeTuplePairsAdmissible.lean:31--47}).

\emph{Historical v8 proof boundary. The current treatment is on page~\pageref{r8:Z4.2.04}.} For each such torsion translate and every multi-index, the Taylor coefficient of the translated canonical vector belongs to \(\overline{\mathbf Q}\). This is the coefficient assertion used for each summand in \nolinkurl{GlobalTheory.TwV.translate_entire_algebraic} (\path{Development/AbelianVarietyTheory/Proofs/SatakeParameters/TranslateEntireAlgebraic.lean:15--18}). The \emph{dictionary} is the premise (B) in this unit and is supplied by \nolinkurl{joint_dict} as above.

% Parent source SHA-256: 726efab7fd295a2dce4fc43b2a38c6ed087511bfb339c03e564ed0607ef7693c
% blueprint/ch_P_r15.tex lines 49 through 121; r15:site-glob
\paragraph{Recovering the global lattice and its adelic span.}\phantomsection\label{r15:site-glob}
Put $\mathcal A=E\otimes_F\mathbb A_{F,f}$ and let $S\subset\mathcal A$
be the subring generated by $\mathcal O_E$ and the finite integral
adeles of $F$. For a lattice $M$, let $\widehat M$ be the $S$-span
of its diagonal image and set $G_g(M)=E^n\cap g\widehat M$ inside
$\mathcal A^n$. The notation $E^n\cap$ uses the diagonal embedding.
This is also the source's \nolinkurl{latAd}: its outer
$\mathcal O_E$-span fixes this submodule. Two scalar facts control the
intersection. First $E\cap S=\mathcal O_E$. To see the nontrivial
inclusion, the trace to $F$ of $xb$ for $x\in E\cap S$ and integral
$b\in E$ is an integral finite adele, hence belongs to $\mathcal O_F$
by the finite valuation criterion. The trace is integral on $S$ since
elements have finite sums of integral $E$ factors times integral $F$
adeles. Testing against an $F$-basis cleared into $\mathcal O_E$ embeds
$E\cap S$ in the $\mathcal O_F$-span of the trace-dual basis. Thus
$E\cap S$ is a finite module over the Noetherian ring $\mathcal O_F$.
It contains $1$ and is stable under multiplication by $x$, so the
determinant trick makes $x$ integral. Second, for
$0\ne d\in\mathcal O_F$, every element of $S$ is congruent to a
diagonal $\mathcal O_E$ element modulo $dS$. For an integral $F$-adele,
choose integral representatives modulo the local powers at the finitely
many primes dividing $d$ and combine them by the Chinese remainder
theorem. At the other primes division by $d$ preserves integrality.
This gives the assertion for the generators of $S$ and hence for $S$.
See \nolinkurl{SDCompan.latAd_eq}, \path{SatakeParameters/IsSDLatLatK.lean:22--24};
\nolinkurl{adelstabT_mem}, \nolinkurl{adelstab_tr_mem}, \nolinkurl{adelstab_dual},
\path{AdelicIwasawaHecke/Adelstab.lean:29--85};
\nolinkurl{csGlob_N_fg}, \nolinkurl{csGlob_int}, \path{AdelicIwasawaHecke/CsGlob.lean:52--101};
\nolinkurl{csGlob_intAdApprox}, \path{AdelicIwasawaHecke/CsGlobApprox.lean:23--86,88--160};
\nolinkurl{csSpan_approx}, \path{AdelicIwasawaHecke/CsSpan.lean:21--48}.

Choose $0\ne b,c_0\in\mathcal O_F$ with
$b\mathcal O_E^n\subset M\subset c_0^{-1}\mathcal O_E^n$:
clear denominators in expressions for the coordinate vectors in the
$E$-span of $M$ and in coordinates of finitely many generators.
Given $x\in E^n\cap\widehat M$, write its diagonal image as a finite
sum of $S$-multiples of these generators. Approximate the scalar
coefficients modulo $bc_0 S$ by integral global coefficients as above.
The resulting $y\in M$ has $x-y\in bS^n\cap E^n=b\mathcal O_E^n$.
Thus $G_1(M)=M$. For general $g$ clear the denominators of its entries
and those of its inverse by one $d\ne0$ in $\mathcal O_F$. Then
\[
 bd S^n\subset g\widehat M\subset (dc_0)^{-1}S^n .
\]
For such a denominator, write each tensor coefficient as a finite sum
of pure tensors. Clear the algebraic denominators in the global factors
and the finitely many nonintegral components in the finite-adele factors.
Multiply the resulting denominators over the entries of $g$ and $g^{-1}$.
The elementary arithmetic used here is the finite-module integrality
criterion, the nondegenerate trace pairing for separable fields and
Dedekind factorization with Chinese remaindering. See J.~S.~Milne,
\emph{Algebraic Number Theory}, v3.01, Chs.~2--3 and 7
(\url{https://www.jmilne.org/math/CourseNotes/ANT301.pdf}).
The intersections with $E^n$ imply that $G_g(M)$ is full and finite:
it contains $bd\mathcal O_E^n$ and is a submodule of a finite
$\mathcal O_E$-module. For $z\in g\widehat M$, approximate each entry
of $dc_0 z$ by an integral global entry modulo $(dc_0)bd S$.
Division by $dc_0$ gives a rational $x$ with $z-x\in bd S^n$.
Thus $x\in G_g(M)$, and each $bd$ times a coordinate basis vector is also
in $G_g(M)$. So $z\in\widehat{G_g(M)}$. The reverse inclusion follows
on generators. We have proved
\[
 \widehat{G_g(M)}=g\widehat M,\qquad G_h(G_g(M))=G_{hg}(M).
\]
The first formula recovers $G_1$ at $g=1$ together with the rational
intersection argument. The second gives inverse transport by $g^{-1}$.
See \nolinkurl{levelaf_exists_denom}, \path{AdelicIwasawaHecke/Levelaf.lean:423--442};
\nolinkurl{csGlob_lattice}, \nolinkurl{csGlob_one},
\path{AdelicIwasawaHecke/CsGlob.lean:26--37,150--173,175--241};
\nolinkurl{csSpan_back}, \nolinkurl{span_recovery_of_integral_adelic_approx},
\path{AdelicIwasawaHecke/CsSpan.lean:50--68,89--115,121--156};
\nolinkurl{Comp.glob_one}, \nolinkurl{Comp.mem_hat_glob}, \nolinkurl{Comp.glob_glob},
\path{GlobalTheory/EqGlobOfHat.lean:35--59}.

% blueprint/ch_P_r15.tex lines 123 through 163; r15:site-dual
\paragraph{Transporting the Hermitian dual and the companion.}\phantomsection\label{r15:site-dual}
We first verify the finiteness hypothesis on duals. In orthonormal
coordinates $c$ take $v_r=c^{-1}e_r$. Since $M$ spans $E^n$,
some nonzero rational integer $n_r$ has $n_r v_r\in M$. Testing a
dual vector $z$ gives $n_r(cz)_r\in\mathcal O_E$. Thus $M^\#$ is
contained in the module generated by $v_r/n_r$ and is finite by
Noetherianity. Clearing one integer $m\ne0$ in the $c$-coordinates
of a generating set of $M$ gives $m\overline{(cx)_r}\in\mathcal O_E$
for every $x\in M$. Conjugation preserves algebraic integers.
Consequently $m v_r\in M^\#$ for every $r$, proving full span.
This argument, including the empty basis when $n=0$, proves
\nolinkurl{OrthDual.isLat_dl} and \nolinkurl{isLat_dl_of_isOrthCo}
(\path{SatakeParameters/IsLatDlOfIsOrthCo.lean:25--57,64--101,103--141}).
The identity $c^*c=J$ gives the required orthonormal coordinates on
expanding $h_{\rm std}(cx,cy)=h_J(x,y)$
(\nolinkurl{isOrthCo_toLE_of_cT_mul_self},
\path{SatakeParameters/IsOrthCoToLEOfCTMulSelf.lean:18--31}).

For $z\in G_g(M^\#)$ and $x\in G_g(M)$ we have
$g^{-1}z\in\widehat{M^\#}$ and $g^{-1}x\in\widehat M$.
Pairing the generators is integral, and the involution preserves $S$:
it preserves $\mathcal O_E$ and fixes the $F$-adele factor.
By sesquilinearity $h_J(g^{-1}x,g^{-1}z)\in S$.
Unitarity identifies this with the diagonal $h_J(x,z)$, so the scalar
intersection just proved puts it in $\mathcal O_E$. This gives
$G_g(M^\#)\subset G_g(M)^\#$. Apply the inclusion to $g^{-1}$ and
$G_g(M)$, then apply $G_g$ and inverse transport, which is valid for
the full dual lattice just proved. The reverse inclusion follows.
Therefore $G_g(M^\#)=G_g(M)^\#$ and self-duality is preserved.
On full lattices $G_g$ is an order isomorphism with inverse $G_{g^{-1}}$,
so it preserves sums. For $a\ne0$ in $\mathcal O_E$,
$G_g(aM)=aG_g(M)$ follows by dividing the rational vector by $a$ inside
the defining intersection and using $\widehat{aM}=a\widehat M$.
It follows that $C_a(G_gM)=G_g(C_aM)$. For $a=0$ both sides reduce
directly to $G_gM$. This accounts for every $a$ in the site theorem.
See \nolinkurl{csDual_sig_mem}, \nolinkurl{csDual_span},
\path{AdelicIwasawaHecke/CsDual.lean:42--110};
\nolinkurl{Comp.glob_sup}, \nolinkurl{Comp.glob_smul}, \nolinkurl{Comp.glob_dl},
\path{GlobalTheory/EqGlobOfHat.lean:69--85,116--138,171--205};
\nolinkurl{Comp.glob_companion}, \nolinkurl{Comp.glob_selfdual},
\path{GlobalTheory/ExistsUnitaryOneDetEq.lean:18--29}.

% blueprint/ch_P_r15.tex lines 165 through 203; r15:site-standard
\paragraph{From the identity chart to the actual kernel lattice.}\phantomsection\label{r15:site-standard}
Set $M_0=C_a(L_b)$. Since $G_1(L_b)=L_b$, the identity-chart hypothesis
says $cM_0=\mathcal O_E^n$. The standard lattice is generated by the
coordinate vectors. Its dual equals itself because their pairings extract
the coordinates, and the pairing of two integral vectors is integral.
Thus $M_0$ is finite, full and $J$-self-dual by pulling back through
the isometry $c$. Explicitly, $y\in M_0^\#$ iff $cy$ pairs integrally
with all of $cM_0$ iff $cy\in cM_0$ iff $y\in M_0$.
Every $G_g(M_0)$ is now a finite, full, self-dual lattice and equals
$C_a(G_gL_b)$. The source finally applies the canonical isometry
$u=\mathrm{stdCo}(J)$, chosen using the existence witnessed by $c$.
For every $z$, membership in $uG_g(M_0)$ is the membership test for
$u^{-1}z$ in $G_g(M_0)$; the isometry identifies their dual tests.
Finite generation and span are preserved under the same linear isomorphism.
By its definition $uG_g(M_0)=\mathrm{latK}(J,a,g,L_b)$. This proves
\nolinkurl{SDCompan.isSDLat_latK} first for any already full self-dual
companion with the stated dual-lattice and coordinate hypotheses, and
then \nolinkurl{KerSelfDual.isSDLat_latK_of_stdL} by the chart hypothesis.
See \nolinkurl{FourthRootStd.stdL_isSD}, \path{SatakeParameters/SpanRangeEmb4EqTop.lean:32--76};
\nolinkurl{OrthSelfDual.isLat_and_dl_eq_of_isSDLat_mvL},
\path{SatakeParameters/IsLatAndDlEqOfIsSDLatMvL.lean:17--26,28--71};
\nolinkurl{SDCompan.isSD_toStd}, \nolinkurl{SDCompan.isSD_core},
\path{SatakeParameters/IsSDLatLatK.lean:33--88,98--100};
\path{SatakeParameters/IsSDLatLatKOfStdL.lean:19--34}.

For a valid tuple at $j\le m$ the site's dictionary supplies the rank
identification with $D.R.N(j)$ and the dependent identification of
$\mathrm{kerN}_j$ with the lattice of its chosen chart and point.
The pinned coordinate satisfies $c^*c=J$, and the joint base law gives
the identity-chart equality above. After transport along the rank
equality the lattice is exactly the site's $\mathrm{kerN}_j$.
The rank is even by \nolinkurl{product_zeroth_first_parts.N_even} (\path{ProductZerothFirstParts.lean:303--305}).
Apply this successively at $j$ and $j+1$ under the single bound $j+1\le m$.
These are the four fields of the admission predicate for the actual pair
in \nolinkurl{pairOK_allSites}.
See \nolinkurl{KerSiteSelfDual.isSDLat_of_heq},
\nolinkurl{KerSiteSelfDual.even_and_isSDLat_kerN},
\path{SatakeParameters/KernelLatticeSelfDualEvenRankAtSite.lean:13--27};
\path{SatakeParameters/KernelSiteLatticeTuplePairsAdmissible.lean:23--31}.

% blueprint/ch_P_r15.tex lines 230 through 292; r15:span-place
\paragraph{Recovering the completed lattice from the simultaneous rational conditions.}
\phantomsection\label{r15:span-place}
Write $\mathcal S$ for the integral tensor subring generated by
$\mathcal O_E$ and the integral finite adeles of $F$, $\widehat L$ for
the $\mathcal S$-span of $L_b$, and $M=\{r\in E^n:x^{-1}r\in\widehat L\}$.
This set is an $\mathcal O_E$-module and agrees with $\mathrm{latAd}(x,L_b)$.
Its $\mathcal S$-span equals $x\widehat L$ as in
Section~\ref{r15:site-glob}. In detail, the bounds for a full finite
lattice and the common denominators of $x,x^{-1}$ give nonzero
$\alpha,\beta\in\mathcal O_F$ with
$\beta\mathcal S^n\subseteq x\widehat L\subseteq\alpha^{-1}\mathcal S^n$.
For $u\in x\widehat L$, lift each coordinate of $\alpha u$ modulo
$\alpha\beta\mathcal S$ to $r_j\in\mathcal O_E$. For an integral
$F$-adele the lifts are obtained at the primes dividing $(\alpha\beta)$ by
approximating in the corresponding completions and applying the Chinese
remainder theorem to those prime powers. At the other primes the modulus
is a unit. The congruence extends to the integral tensor ring because
its two sets of generators have lifts and congruence is preserved by
addition, subtraction and multiplication. Then
$u=\alpha^{-1}r+\beta w$, $w\in\mathcal S^n$. The rational vector
$\alpha^{-1}r$ belongs to $M$ by the lower bound, as does each
$\beta e_j$. Hence all summands lie in the $\mathcal S$-span of $M$.
The reverse inclusion follows from the defining condition on $M$.
See \nolinkurl{csApprox_global}, \nolinkurl{integral_adele_approximation_holds}
(\path{AdelicIwasawaHecke/CsApprox.lean:20--109}), \nolinkurl{csSpan_approx},
\nolinkurl{csSpan_back}, \nolinkurl{span_recovery_of_integral_adelic_approx}
(\path{AdelicIwasawaHecke/CsSpan.lean:21--63,64--156}), and
\nolinkurl{latAd} (\path{OrthogonalChartLevels/IsOrthCo.lean:17--21}).

Write $R_v=E\otimes_F F_v$, $S_v=\mathrm{intV}(E,v)$, $x_v$ for
the matrix over $R_v$ and $L_{b,v}=S_v L_b$. Then
$S_vM=x_v L_{b,v}$. For the inclusion from the right,
take $z=x_v h$, $h\in L_{b,v}$. Extend the coefficients of a local
linear combination for $h$ by zero off $v$ to form an adelic vector in
$\widehat L$. On an embedded $E$-generator this extension multiplies it by the place
idempotent, and on $1\otimes b$ it extends the integral local coefficient
$b$ by zero. Sums and products stay in the integral adelic ring. Multiply the
extension by $x$, apply the global span identity, and project at $v$.
Its projection is $z$. Conversely each global generator of $M$ maps into
$x_v L_{b,v}$ by applying $x_v^{-1}$, and this condition is $S_v$-linear.
These are \nolinkurl{LocCarry.mem_latS_latAd_iff}
(\path{SatakeParameters/MemLatSLatAdIff.lean:37--72}) and
\nolinkurl{IE_latS}, \nolinkurl{prV_latS} (\path{GeometryDefinitions/PNeatB.lean:24--72}).

Project $R_v$ onto its first split factor and transport the result to
the field $Q_v$ of the chart. Call this map $f_v$, so its restriction
to $E$ is $\rho_v$. Its image on $S_v$ is exactly $W_v$ (identified
with its image in $Q_v$). Indeed $\mathrm{SplAlg}$ sends $1\otimes z$
to $(z,z)$. The other generators $a\otimes1$, $a\in\mathcal O_E$,
have integral images by their monic equations over $\mathbb Z$ and
integral closedness of the local valuation ring. Every local integer
is the image of $1\otimes z$ for that integer $z$, and $\mathrm{LocInt}$
makes the resulting ring precisely the chart's $W_v$. It follows
in both directions, by lifting coefficients of finite sums, that
$f_v(S_vL)=\operatorname{span}_{W_v}\rho_v(L)$ for any lattice $L$.
Since $f_v(x_vh)=f_v(x_v)f_v(h)$ entrywise, the preceding equality gives
\[
 \operatorname{locSp}_v(M)=f_v(x_v)\operatorname{locSp}_v(L_b).
\]
See \nolinkurl{ph_tensInt_onto}
(\path{SatakeParameters/TensorFirstFactorIntegersOntoLocalIntegers.lean:14--49}),
\nolinkurl{locSp_eq_image} (\path{SatakeParameters/LocSpEqImage.lean:16--66})
and \nolinkurl{latAdLoc_of} (\path{SatakeParameters/LatAdLocOf.lean:17--60}).

% blueprint/ch_P_r15.tex lines 294 through 326; r15:span-frame
\paragraph{The companion, the chart and the order of multiplication.}
\phantomsection\label{r15:span-frame}
Transport of the Hermitian dual gives $M^\#=\mathrm{glob}_x(L_b^\#)$
for unitary $x$ (Section~\ref{r15:site-dual}). Scaling commutes with
$x^{-1}$, so the level condition $T L_b^\#\subseteq L_b$ gives
$T M^\#\subseteq M$. For the companion $M_a=M+(a)M^\#$ this yields
$M\subseteq M_a$ and $T M_a\subseteq M$ for every integral $a$.
Each $\rho_v(z)$ with $z\in M_a$ therefore lies in
$\operatorname{locSp}_v(M)$ after multiplying $\rho_v(Tz)$ by the
inverse of $T$ in $W_v$. The two local spans agree. See
\nolinkurl{smul_dualSet_latAd_mem}
(\path{SatakeParameters/SmulDualSetLatAdMem.lean:17--55}),
\nolinkurl{CompSp.locSp_compan_eq}
(\path{SatakeParameters/LocSpCompanEq.lean:17--59}),
and \nolinkurl{compan}, \nolinkurl{stdCo}, \nolinkurl{toStd}
(\path{OrthogonalChartLevels/IsOrthCo.lean:33--43}).

By $\mathrm{LbAtB}$ the columns of $B_v=\rho_v(B)$ span the first
factor of $L_b$. Projecting the place clause of $\mathrm{IsPt}$ gives
$f_v(x_v)=B_vg_vB_v^{-1}$. The chosen point therefore gives
$\operatorname{locSp}_v(M_a)=\operatorname{colSpan}_{W_v}(B_v g_v)$.
Finally $\mathrm{latK}$ is $M_a$ sent through the standard orthogonal
coordinate map with matrix $C$. On each rational generator $z$,
$\rho_v(Cz)=\rho_v(C)\rho_v(z)$, and the $W_v$-linear image of the
span of the columns of a matrix $Y$ is the span of the columns of its
left multiple. This proves the claimed frame formula, with no assumption
that $a$ is a unit. It applies to every $v$ at which $T$ is a unit.
See \nolinkurl{FacQPt.facQ_isPt} (\path{SatakeParameters/FacQIsPt.lean:19--33}),
\nolinkurl{SpanPt.colSpan_map}, \nolinkurl{SpanPt.locSp_latAd_isPt}
(\path{SatakeParameters/LocSpLatAdIsPt.lean:17--29}),
\nolinkurl{locSp_map} (\path{SatakeParameters/MapApplyEqMulVec.lean:9--24}),
and \nolinkurl{KerSp.locSp_latK_isPt}
(\path{SatakeParameters/KernelLatticeFirstFactorLatticeTuplePoint.lean:21--30}).

% blueprint/ch_P_r15.tex lines 328 through 386; r15:span-steps
\paragraph{The zero and successor branches on the same tuple.}
\phantomsection\label{r15:span-steps}
Work with the tied entry $t$, the served prime $p$, the specified
standard local data, the joint datum and its dictionary. For each valid
tuple and place, \nolinkurl{Pend.latAt_isLattice}
(\path{SubstZOfXCoordZA.lean:289--293}) gives a matrix with unit
determinant whose columns span the tuple lattice. Hence this matrix is invertible over the local
field and the basis matrix selected in $\mathrm{basesAt}$ has that same
span. The chosen representative $J.\mathrm{pt}$ actually represents this
component and these bases: the goodness record gives such a unitary
point and the defining choice reads its specification. Also
$\mathrm{GIdx}$ supplies $p\nmid T$ and primality of $p$
(\path{LocalRingCardinality/RhoDetPrimeEqMuTildeMulChiPow.lean:68--69}).
This gives $\gcd(T,p)=1$ and $T>0$. Only this coprimality projection
of \nolinkurl{gidx_facts} enters the span argument. Its inertness
projection reaches \nolinkurl{Served.served_facts}, accounting for 99
weighted lines unused here (\path{SatakeParameters/GidxFacts.lean:16--24}).
In the local ring the image of $p$ is a nonunit, so a B\'ezout relation
forces the image of $T$ to be a unit. If $Q^*Q=J$ is the pinned
coordinate identity, expansion of the Hermitian form shows $Q$ gives
orthogonal coordinates. See \nolinkurl{Joint.basesAt_span}
(\path{SatakeParameters/LatticeTupleBasesSpanAtPlaces.lean:10--16}),
\nolinkurl{exists_unitary_pt} (\path{SatakeParameters/ExistsUnitaryPt.lean:14--16}),
\nolinkurl{XdInert.gidx_facts} (\path{SatakeParameters/GidxFacts.lean:16--24}),
\nolinkurl{natCast_isUnit_localRing}
(\path{SatakeParameters/IsUnitNatCastOfCoprimeLocal.lean:16--26}),
and \nolinkurl{isOrthCo_toLE_of_cT_mul_self}
(\path{SatakeParameters/IsOrthCoToLEOfCTMulSelf.lean:18--31}).

At zero, rewrite the dictionary's equality of the actual kernel lattice
using $J.\mathrm{lev0}$, and apply the preceding formula to the
level-zero chart with $J.\mathrm{alg0}$, $J.\mathrm{one.int}$,
$J.\mathrm{lb0}$ and the point just described. At step $j+1$ one has
$j+1\le D.R.m$. Decompose the tuple's component by the chart equivalence
$e_C^{-1}$. The transported chart equality for this component identifies
both its $T$ with the served level and its base lattice with the lattice
in the dictionary. Use the component's own
$J.\mathrm{alg}$, $H.\mathrm{core.int}$ and $J.\mathrm{lbB}$ and its
$\mathrm{IsPt}$ for the actual successor point. In both cases the
dictionary supplies the rank equality and the corresponding equality
of lattices after reindexing. Transport by that finite-index bijection
preserves generators. Thus the right side becomes the image of
$\mathrm{latAt}(v,i)$ under the frame. At the successor the valid branch
of $\mathrm{chAt}$ identifies its Gram matrix, $B$, and local embedding
with those of the selected component. The equality $J.\mathrm{iw}$
identifies that component embedding with $J.\rho(v)$, giving exactly
$\mathrm{frM}(j+1,\mathrm{La},v)$ in this formula. These arguments take
the given tuple and place in both adjacent positions. See
\nolinkurl{Joint.isPt_pt_zero}, \nolinkurl{Joint.pt_succ_isPt}
(\path{SatakeParameters/LatticeTupleAdelicPointStepZero.lean:12--35},
\path{SatakeParameters/LatticeTupleAdelicPointStepSucc.lean:14--42}),
\nolinkurl{T.lattice_tuple_chart_succ_eq_component_chart}
(\path{SatakeParameters/LatticeTupleChartSuccEqComponentChartStatement.lean:9--10})
and its implementation \nolinkurl{chart_succ_eq}
(\path{SatakeParameters/LatticeTupleChartSuccEqComponentChart.lean:22--27}),
\nolinkurl{locSp_reindex} (\path{SatakeParameters/ColSpanMapMulVecLin.lean:15--26}),
and \nolinkurl{Joint.locSp_kerN_zero}, \nolinkurl{Joint.locSp_kerN_succ}
(\path{SatakeParameters/FrameLocalLatticeSpanKernelLatticeZero.lean:19--34},
\path{SatakeParameters/FrameLocalLatticeSpanKernelLatticeSucc.lean:18--54}).

% blueprint/ch_P_r15.tex lines 388 through 413; r15:datum
\paragraph{The moved datum has the required period congruence.}
\phantomsection\label{r15:datum}
At $j<\mathrm{mm}=D.R.m$, fix both valid adjacent tuples, their points,
and their actual two kernel lattices $N,N'$. The dictionary at the two
indices gives rank identifications $h_1,h_2$ and identifies both lattices
with the chart lattices. Put $k=(t.\mathrm{rd}(u,\mu)).\mathrm{lev}$ and
$S_E=\mathrm{ssRec}(t,u,\mu)$. The claim concerns two linear maps $z,z'$:
if at every $P\in S_E$ the difference
$k^{-1}(z-z')$ belongs to the localization of
$\mathrm{perO}(\mathrm{aaOf}(D,t.cc),N,N')$ at $P$, their moved
datum values agree. In $\mathrm{xiM}$ the positive branch of
$\mathrm{xi0}$ is selected by $j+1\le D.R.m,h_1,h_2$.
The two frame changes are inverses on the source and target of the
linear map. Composing them cancels both, leaving the site's adelic
datum on these charts and points evaluated at $\mathrm{homRe}(h_1,h_2,z)$.
The site's $S.\mathrm{law.loc}$ already imposes the displayed congruence
there. Identify the two ranks, then substitute the two dictionary lattice
equalities. The localization conditions are now identical, and
$\mathrm{homRe}(\mathrm{refl},\mathrm{refl},z)=z$. This proves
$\mathrm{DatLoc}$ on the original pair with the original level.
See \nolinkurl{DatLoc} (\path{GaloisGeometry/NatCastMemOfPlace.lean:14--18}),
\nolinkurl{inFr_symm_inFr_cancel}, \nolinkurl{moved_pair_datum_eq_adelic_datum_alt}
(\path{SatakeParameters/MovedPairDatumEqAdelicDatum.lean:13--24}),
\nolinkurl{datLoc_homRe} (\path{SatakeParameters/DatLocHomRe.lean:14--23}),
and \nolinkurl{TwU.xiM_datLoc}
(\path{SatakeParameters/MovedPairDatumDependsOnlyOnLevel.lean:13--28}).

% blueprint/ch_P_r15.tex lines 415 through 644; r15:translated
\subsection{Z4.2.04. Translated coefficients of the natural vector}
\label{r15:translated}
Fix the curve ideal $\mathfrak c$, its discriminant root $\delta$, a base $B$,
an $\mathrm{IsChi}$ character $\chi$, and an admitted pair $(L,L')$. If
$m x\in A_{L,L'}=\mathrm{La}(L,L')$ for an integer $m>0$, then
$U_x(B.e^{\rm nat}_\chi(L,L'))$ is entire and every normalized Taylor
coefficient at zero belongs to $k=\overline{\mathbf Q}$.
This is \nolinkurl{TwV.translate_entire_algebraic}
(\path{SatakeParameters/TranslateEntireAlgebraic.lean:15--18}).

\emph{Hand proof and hypotheses.} We translate all four variables in
the defining tensor identity. Lean first proves the unshifted statement
and then calls \nolinkurl{AlgTranslate.coefIn_translate}. For the elliptic
translate we use a division-quotient recurrence in place of
\nolinkurl{ThetaConj.cmTheta_taylorCoeff_algebraic_and_conjugate}
(\path{SatakeParameters/CoefInTranslate.lean:53--85},
\path{VacuumEntireAndNormalizedJetsAlgebraic.lean:311--316},
\path{CmThetaTaylorCoeffAlgebraicAndConjugate.lean:263--290,379--389}).
The selected function $\theta=B.\mathrm{thMM}$ is an epsilon choice.
Given $B.\mathrm{StdOK}$ and $B.\kappa\ne0$, the coefficient branch
obtains its four-copy identity from \nolinkurl{WeilVac.thMM_spec},
which selects the witness supplied by \nolinkurl{ThetaVec.thetaVector}.
This second implementation of the proposition in Section~\ref{r16:ktheta}
adds 1,389 core lines, including 1,348 theorem lines, to the round 16
registered union in the residual Pair cone. The parallel selection
\nolinkurl{Base.thMM_spec} via \nolinkurl{Ker.kThetaVector} enters the
Pair cone through the entireness clause. Section~\ref{r16:stdok}
supplies the standard data in both applications
(\nolinkurl{thM}, \nolinkurl{Base.thMM}, \path{KernelLattices/LateralTheory.lean:23--29,155},
\path{IwasawaDecomposition/PeriodLatticeSesquilinearFormIsHerm.lean:437--498},
\path{IwasawaDecomposition/ExistsWeilOThZeroTriv.lean:33--36},
\path{VacuumEntireAndNormalizedJetsAlgebraic.lean:197--200},
\path{KernelLocalFactors/IsAdmissibleFlagPairOfUnitriangularLeftAutomorph.lean:27--29,75--81,218--220}).

There are two further differences in the translation argument. Lean
gets entireness from \nolinkurl{vector_bound_link}. The finite sums of
translates below remain entire term by term
(\path{GlobalTheory/GL3/ProvisionExtensionLink.lean:38--39},
\path{SatakeParameters/TranslateEntireAlgebraic.lean:15--18}).
For rational phases Lean uses \nolinkurl{RationalPoints.rational_points},
which takes two integrality clauses from the seven-part
\nolinkurl{ThetaLine.theta_line}. Here only the alternating pairing on
the standard lattice is needed, supplied by \nolinkurl{Base.latAll}
(\path{DivisorProofs/RationalPoints.lean:18--43},
\path{FoundationProofs/ThetaSpaceSelfDualPairAtMostLine.lean:276--294},
\path{KernelLocalFactors/IsAdmissibleFlagPairOfUnitriangularLeftAutomorph.lean:623--644}).
In review A's ordered partition of the translation cone, after setting
aside the original registered union and the two algebraicity detours
above, the entireness branch accounts for 270 core lines, including
247 theorem lines. Of the remaining 6,237 core lines, the
\nolinkurl{theta_line} detour accounts for 2,046, including 2,015
theorem lines (\path{out/R16_review_A.md}, Section~5.1).

\paragraph{The elliptic germ.}
Write $\Lambda=\mathrm{lamOf}(D,\mathfrak c,\delta)$ and $f=\mathrm{sgStar}(\Lambda)$.
The epsilon choice \nolinkurl{TaylorLogMetric.sgStar} equals
$\sigma_\Lambda(z)\exp(-s_2z^2/2)$. This function has the required
multiplier, oddness and derivative one at zero. The one-pole
uniqueness argument in Section~\ref{r16:ktheta-line} determines the
normalized function.
Lean identifies these functions in \nolinkurl{OddTheta.sgStar_eq} and
their normalized half-period translates in \nolinkurl{thetaZero_periodLattice_eq}
(\path{IsDiscriminantTwelfthRootOfIdealLattice.lean:238--240},
\path{OddThetaFunctionExistenceUniqueness.lean:97--105},
\path{VacuumEntireAndNormalizedJetsAlgebraic.lean:187--189}).
The period datum has discriminant $g_2^3-27g_3^2=1$ and is stable under
$\mathcal O_K$. Its lattice ratio is in the upper half-plane in the
imaginary quadratic field $K$ (a ratio of two generators of the CM ideal,
unaffected by scaling). Milne, \emph{Modular Functions and Modular Forms},
v1.31 (2017-03-22), Corollary~12.16 makes its $j$ an algebraic integer. Thus
$j=1728g_2^3$ gives $g_2,g_3\in k$. For $u\in\Lambda_{\mathbf Q}\setminus\Lambda$,
uniformization sends $u$ to the point $(\wp(u),\wp'(u))$ on
$y^2=4x^3-g_2x-g_3$. Multiplication by an integer killing $u$ has
finite kernel defined over $k$. Both coordinates lie in $k$, which
is algebraically closed. See Silverman, \emph{The Arithmetic of Elliptic
Curves}, 2nd ed., III.6.4 and VI.3.6, and Milne's corollary
(\url{https://www.jmilne.org/math/CourseNotes/MF.pdf}).
The corresponding Lean results are \nolinkurl{CurveLattice.lattice_of_curve}
and \nolinkurl{CurveAlgebraic.integral_j_and_algebraic_data}
(\path{IdealPeriodLatticeOfCurveDatum.lean:136--157},
\path{IntegralJAndAlgebraicData.lean:35--109,126--150}).

The classical sigma has simple zeros at $\Lambda$, logarithmic derivative
$\zeta$ and $-\zeta'=\wp$. Multiplying sigma by a quadratic exponential
to obtain the Hermitian-normalized $f$ preserves its addition identity
and zeros. The derivative is normalized to $f'(0)=1$ and $f$ is odd.
The logarithmic second derivative has the form
$(\log f)''=-(\wp+s)$ for a constant $s$.
For the classical laws, the nonsingular lattice discriminant and the
sigma addition formula see NIST DLMF \S\S23.2(ii)--(iii), 23.3(i),
23.10.3 (\url{https://dlmf.nist.gov/23.2},
\url{https://dlmf.nist.gov/23.3}, \url{https://dlmf.nist.gov/23.10.E3}).
The multiplier in the project's normalization is the one in
\nolinkurl{OddTheta.sgStar_laws} (\path{OddThetaFunctionExistenceUniqueness.lean:115--147}).

The normalized Laurent expansion of $\wp$ lies in $k((z))$.
The terms following $z^{-2}$ start with $g_2z^2/20+g_3z^4/28$,
and $\wp''=6\wp^2-g_2/2$ determines the subsequent terms.
For $\rho\in\mathrm{Aut}(\mathbf C/k)$ apply $\rho$ to the
formal coefficients. The logarithmic derivative identity, oddness and
the coefficient of $z$ give $f^\rho=f\exp(cz^2)$ with
$2c=s-\rho(s)$. Choose the host CM scalar $\alpha$ with
$\alpha^2=-d$, $|\alpha|^2=d>0$, $\alpha\Lambda\subseteq\Lambda$.
By the multiplier identity $f(\alpha z)^2/f(z)^{2d}$ is even and elliptic.
Its poles all lie in $\Lambda$, so it equals $P(\wp)$ for a polynomial:
subtract powers of $\wp$ to remove the even principal part, leaving a
holomorphic elliptic function. Conjugating the formal identity gives
\[
 \exp(2c(\alpha^2-d)z^2)=P^\rho(\wp(z))/P(\wp(z)).
\]
The right side is meromorphic and periodic. Its identity with the
entire left side extends from the germ. Periodicity of the logarithmic
derivative forces $c=0$ since $\alpha^2-d\ne0$. Thus every coefficient
of $f$ belongs to the fixed field $\mathbf C^{\mathrm{Aut}(\mathbf C/k)}=k$
(Milne, \emph{Fields and Galois Theory}, v5.00, Thm.~9.29, since $k$ is
perfect and $\mathbf C$ algebraically closed;
\url{https://www.jmilne.org/math/CourseNotes/FT500.pdf}).
See \nolinkurl{conj_sgStar}, \nolinkurl{exists_poly}, \nolinkurl{twist_eq_zero},
\nolinkurl{galSg_all} (\path{GaloisTaylorSeriesCanonicalOddThetaAll.lean:389--433,441--515,555--649})
and \nolinkurl{complex_mem_range_of_forall_algEquiv}
(\path{ExistsRingEquivOfIsTranscendenceBasis.lean:107--110}).

\paragraph{Torsion translates.}
For $u\in\Lambda_{\mathbf Q}\setminus\Lambda$, set
$\Psi_j=f(j(z+u))/f(z+u)^{j^2}$. Then $\Psi_1=1$ and
$\Psi_2=-\wp'(z+u)$. The torsion coordinates and the differential
equation for $\wp$ put both germs in $k((z))$. The sigma addition
formula and the chain rule give
\[
 (j+1)^2\Psi_j\Psi_{j+2}
    =\Psi_{j+1}\Psi_{j+1}''-(\Psi_{j+1}')^2\qquad (j\ge1).
\]
Each $\Psi_j$ is a nonzero germ, so division in $k((z))$ determines
the next one even when the value at zero vanishes. Choose $m>0$ with
$mu\in\Lambda$. For $h=2m$ the sign at $hu$ is $+1$ since
$hu/2\in\Lambda$. The definition of translation then gives
$(U_uf)^{h^2}\Psi_h=U_{hu}f(hz)=f(hz)$. If
$U_uf=\sum c_jz^j$ and $p=h^2$, division by $\Psi_h$ yields
$c_0^p\in k$, and $c_0\ne0$ since $f(u)\ne0$. The degree-$j$ coefficient
contains $pc_0^{p-1}c_j$ plus a polynomial in earlier coefficients.
Induction gives all $c_j\in k$. For $u\in\Lambda$ the multiplier
already gives $U_uf=\pm f$. For the even half-period theta
$e_0=U_{t_0}f/(U_{t_0}f)(0)$, $t_0=\omega_0/2\notin\Lambda$,
the denominator is nonzero algebraic. Composing $U_u$ and $U_{t_0}$
introduces a root of unity because positive multiples of both shifts
are periods and the alternating form is integral there.
See \nolinkurl{IsSigmaLike.divQuot_rec}, \nolinkurl{IsSigmaLike.divQuot_ne}
(\path{SecondDerivSigmaProductAtZero.lean:127--168}),
\nolinkurl{thetaTranslate_canonicalOddTheta_of_half_mem}, \nolinkurl{thetaTranslate_natCast_mul}
(\path{GaloisConjTaylorThetaTranslate.lean:42--49,51--64}),
\nolinkurl{phase_mem} and \nolinkurl{thetaTranslate_evenTheta_eq_sigmaStar_translate}
(\path{AlgebraicTaylorTorsionTranslateEvenTheta.lean:15--26,40--63,73--77}).

\paragraph{Four independent translated variables.}
Put $L_M=B.\mathrm{lamM}$, $\gamma=B.g_C$ and let $a\in(L_M)_{\mathbf Q}$.
In the identity $\theta^{\otimes4}=\kappa K$ supplied by
\nolinkurl{WeilVac.thMM_spec}, the kernel $K$ is the normalized finite
sum of the products $U_b\prod e_0$ composed with $\gamma$,
with signs indexed by the relative quotient of the two period lattices.
The signs, index square root and $\kappa$ are algebraic: the powers
defining $\kappa$ are ratios of algebraic conjugates. Translate all
four variables by $\mathbf a=(a,a,a,a)$. Orthogonality turns the left
side into $(U_a\theta)^{\otimes4}$ and each summand into the product
translated first by $b$, then by $\gamma\mathbf a$, composed with
$\gamma$. The denominator inclusion for $\gamma L_M^4$ puts both
shifts in the rational span of the product periods. The resulting
phase is a root of unity. The complex matrix of $\gamma$ has algebraic
entries, since it takes a basis of algebraic CM embedding vectors to
algebraic vectors. Algebraic germs stay algebraic under this linear
substitution. At multi-index $(\beta,\beta,\beta,\beta)$ there is
exactly one coefficient from each variable block, giving
$([z^\beta]U_a\theta)^4\in k$. All its fourth roots lie in $k$.

For the standard pair $S,S'$, \nolinkurl{Base.lat0OK} gives
$A_{S,S'}=\mathrm{powLat}(L_M,n,m)$: testing against coordinate
vectors puts each period entry in $L_M$, while single-entry maps and
summing over matrix units give the reverse inclusion. The same
positive multiplier therefore puts each entry of a rational-saturation
shift in $L_M$. The Hermitian form is a sum over the disjoint
matrix entries, hence translating the product defining $B.\mathrm{thZero}$
translates each factor separately. Its jets are in $k$, including
the empty product. See \nolinkurl{Base.lat0OK}
(\path{KernelLocalFactors/IsAdmissibleFlagPairOfUnitriangularLeftAutomorph.lean:369--439}),
\nolinkurl{WeilVac.thMM_spec} (\path{IwasawaDecomposition/ExistsWeilOThZeroTriv.lean:33--36}),
\nolinkurl{FourthRootStd.up} (\path{SatakeParameters/SpanRangeEmb4EqTop.lean:613--652}),
\nolinkurl{alg_of_span}, \nolinkurl{kappa_mem} and the unshifted Lean
results \nolinkurl{thetaVector_fourthPower_algebraicGerm}, \nolinkurl{thZero_coefIn}
(\path{VacuumEntireAndNormalizedJetsAlgebraic.lean:144--160,170--177,197--258}).

\paragraph{The two intertwiners and the pair character.}
Let $A_{0,\mathbf Q}$ be the saturation of the standard period lattice.
Each intermediate period lattice is contained in $A_{0,\mathbf Q}$.
Indeed the induced CM embeddings express each period generator as
$\delta\,\mathrm{embV}(M)$ for a rational matrix $M$. On the standard
lattices the trace-dual family for the half-different spans each matrix
entry over $K$. Clear denominators in the finite expansion of $M$.
The norm $|\mathcal O_K/\mathfrak c|$ is nonzero and belongs to
$\mathfrak c$, so multiplying by it as well puts the resulting generator
in $A_{S,S'}$. This proves the inclusion for all sums of generators.
See \nolinkurl{exists_dual_family_hom} (\path{ExistsTraceDualFamilyHomLattice.lean:64--78}),
\nolinkurl{klat_rat}, \nolinkurl{klat_mul_mem} (\path{FoundationProofs/ThetaTranslationIntertwinerUnitary.lean:94--158}),
\nolinkurl{RationalPoints.rational_points} (\path{DivisorProofs/RationalPoints.lean:18--43})
and \nolinkurl{la_le_sat} (\path{VacuumEntireAndNormalizedJetsAlgebraic.lean:290--300}).

At the admitted standard pair, \nolinkurl{Base.latAll} supplies
$E_H(A_0,A_0)\subseteq\mathbf Z$ (Section~\ref{r14:lat-integral},
\path{KernelLocalFactors/IsAdmissibleFlagPairOfUnitriangularLeftAutomorph.lean:623--634}).
For any two shifts in $A_{0,\mathbf Q}$, clearing their positive
denominators therefore makes the alternating pairing integral, so the phase
in their Heisenberg product is a root of unity. Entire functions
with algebraic jets at every such translate are thus closed under
these translations and finite algebraic linear combinations.
The intertwiner averages translates indexed by the receiving
lattice modulo its intersection with the sending lattice, with
coefficients $\pm[ A_2:A_1\cap A_2]^{-1/2}$; the source defines
the zero-index case as zero. Every representative is in the saturation
by the preceding inclusion. Apply this twice to $B.\mathrm{thZero}$
to obtain $B.\mathrm{thMoved}$. Its extra scalar
$\mathrm{chiFac}$ is algebraic: at each admitted lattice the relative
determinant ideal has norm one by self-duality, and a positive
class-number power and the principal-ideal law for $\chi$ give a
nonzero algebraic power of its character value (Section~\ref{r8:Z4.2.04}).
Finally $mx\in A_{L,L'}$
implies $x\in A_{0,\mathbf Q}$ by multiplying once more. Translation
of an entire function multiplies its affine translate by an entire
exponential. This proves both conclusions for the same given pair.
See \nolinkurl{TA.comb}, \nolinkurl{TA.transl}, \nolinkurl{TA.iop},
\nolinkurl{ta_thMoved} (\path{VacuumEntireAndNormalizedJetsAlgebraic.lean:262--316}),
\nolinkurl{chiFac_mem} (\path{TranslateMomentsAllOfJetMomR.lean:324--339}),
and \nolinkurl{ThLoc.Alg.algTAll} (\path{GaloisGeometry/AlgTAll.lean:12--26}).

% blueprint/ch_P_r15.tex lines 646 through 720; r15:pair
\subsection{The completed Z4.2.04 adjacent-site input}\label{r15:pair}
Fix the tied datum, dictionary, site and valid tuples at the adjacent stages
$j,j+1$, with $j+1\le m$. At both steps,
Section~\ref{r15:site} gives the self-dual lattice and the site rank
formula gives even rank. Section~\ref{r15:datum} gives congruence of
the moved datum. At each finite place, Section~\ref{r15:span-steps}
identifies the localized lattice for each tuple with the image under
its own frame, including the separate formula at step zero. Each
frame is a product of invertible local coordinate maps and therefore
has unit determinant in the local field $Q_v$
(\nolinkurl{isUnit_det_frameMatrix_ofCharts},
\path{SatakeParameters/IsUnitDetFrameMatrixOfCharts.lean:14--24}).
These statements refer to the same two tuples.
The tied level is positive and all its prime divisors belong to the specified
bad set by \nolinkurl{sNum_lev_site}. The site's integrality law applies to
its adelic datum at every argument, hence also after the frame changes in
the moved datum (\nolinkurl{Joint.site_base_datum_integral_off_bad_primes},
\path{SatakeParameters/SiteFiniteDataLevelPosBadPrimeFactors.lean:13--15},
\path{SatakeParameters/SiteBaseDatumIntegralOffBadPrimes.lean:13--18}).
Together with admission these give the hypotheses of
\nolinkurl{kerTh_entire_algCoef}
(\path{SatakeParameters/KerThEntireAlgCoef.lean:22--40}).
The complex kernel is entire and has algebraic Taylor coefficients
by the finite Fourier expansion of Section~\ref{r14:kernel-finite}
and the translated coefficient check in Section~\ref{r15:translated}.
The statement wrappers in the Pair proof can be read against these
implementations as follows. \nolinkurl{T.kernelLattice_selfDual_even_rank_at_site}
exports \nolinkurl{KerSiteSelfDual.even_and_isSDLat_kerN}
(\path{KernelLatticeSelfDualEvenRankAtSiteStatement.lean:9--10});
\nolinkurl{T.frame_localLattice_span_kernelLattice_zero} and
\nolinkurl{T.frame_localLattice_span_kernelLattice_succ} export
\nolinkurl{Joint.locSp_kerN_zero} and \nolinkurl{Joint.locSp_kerN_succ}
(\path{FrameLocalLatticeSpanKernelLatticeZeroStatement.lean:9--10},
\path{FrameLocalLatticeSpanKernelLatticeSuccStatement.lean:9--10}).
Inside those two calls,
\nolinkurl{T.kernelLatticeFirstFactor_latticeTuplePoint} and
\nolinkurl{T.lattice_tuple_bases_span_at_places} export
\nolinkurl{KerSp.locSp_latK_isPt} and \nolinkurl{Joint.basesAt_span}
(\path{KernelLatticeFirstFactorLatticeTuplePointStatement.lean:14--15},
\path{LatticeTupleBasesSpanAtPlacesStatement.lean:9--10}).
\nolinkurl{T.movedPairDatum_congr} exports \nolinkurl{TwU.xiM_datLoc}
(\path{MovedPairDatumDependsOnlyOnReaderLevelStatement.lean:9--10}).

The generic sign test for the quadratic diagonal and the Osgood
regularity step appear at Sections~\ref{r14:kernel-lattice}
and~\ref{r14:kernel-jets}. The pair scalar is algebraic and the character
factor is algebraic and nonzero by the local determinant calculation in
Section~\ref{r8:Z4.2.04} (\nolinkurl{latticeTuplePairScalar_mem_algebraicClosure_alt},
\path{SatakeParameters/LatticeTuplePairScalarMemAlgebraicClosure.lean:13--26}).
Indeed \nolinkurl{bwB_gramN_det} gives $\det J\ne0$ by
nondegeneracy. Taking determinants in $g^*Jg=J$ and cancelling
$\det J$ gives $\overline{\det g}\det g=1$, so the site's algebraicity
law applies to both point determinants in $\mathrm{scal0}$
(\path{PBwLocBases.lean:37--58}).
Taking these fields in order gives
$J.\mathrm{Dict}(\mathrm{std})\Longrightarrow S.\mathrm{Pair}$ for
\nolinkurl{GlobalTheory.pairOK_allSites}
(\path{SatakeParameters/KernelSiteLatticeTuplePairsAdmissible.lean:23--47};
\nolinkurl{PairOK}, \path{GaloisGeometry/NatCastMemOfPlace.lean:56--69}).
Removing the cones of Z4.2.01, .03 and .05 leaves 1,508 core
lines for site self-duality and 1,885 for the local span. Their union
has 2,136 core lines, including 2,004 theorem lines.
R15 review A confirmed all 26 site declarations. R16 review A checked
the two remaining inputs in Sections~\ref{r16:ktheta} and~\ref{r16:stdok}
and their application in Section~\ref{r15:translated}. After its
corrections E1--E7, Z4.2.04 is complete without an open condition.
In the residual Pair cone the union below \nolinkurl{Ker.kThetaVector},
\nolinkurl{ThetaVec.thetaVector}, \nolinkurl{FourthRootStd.stdOK}
and \nolinkurl{SO.stdOK} has 8,373 core lines, including 7,887
theorem lines. Including the two three-line selection wrappers
\nolinkurl{Ker.thM_spec} and \nolinkurl{Ker.isThM_ne_zero} gives
8,379 and 7,893 respectively. Cutting the two wrappers, the two
standard-data heads and \nolinkurl{ThetaVec.thetaVector} from the
Pair traversal removes 3,654 residual core lines, including 3,477
theorem lines, that have no other path from the Pair root.

% Parent source SHA-256: 4eafc9a56042cefeaca6da8827cc9f3401b6e853ea24ff22b978e0b558a543dd
% blueprint/ch_P_r16.tex lines 12 through 259; r16:stdok, r16:std-char, r16:std-isom, r16:std-dual, r16:std-indices
\subsection{The standard data on a given base}
\label{r16:stdok}
Fix a host $D$, a curve ideal $\mathfrak c$, a discriminant twelfth root
$\delta$ for that ideal, and \emph{any} $B:\mathrm{Base}(D,\mathfrak c,\delta)$.
Both \nolinkurl{FourthRootStd.stdOK} and \nolinkurl{SO.stdOK} assert
$B.\mathrm{StdOK}$ for precisely these data
(\path{SatakeParameters/SpanRangeEmb4EqTop.lean:700--701},
\path{GlobalTheory/GL3/PStdOK.lean:338--339}). Write $P$ for the lattice
of $B.L$, equal by $B.\mathrm{lat}$ to $\delta\phi_0(\mathfrak c)$, and
$I=\mathfrak b^{-1}=D.\mathrm{host}.\mathrm{AA0}$. The one-block lattice
$M$ is the additive group generated by $p(\tau y)_\tau$, $p\in P$,
$y\in I$. The embeddings $\tau:E\hookrightarrow\mathbf C$ extend
the specified embedding of $K$. In four blocks set
\[
 A=P^{\mathrm{Fin}(4)\times\mathrm{Tau}},\qquad
 A'=g(M^4),\qquad g=B.g_C.
\]
Let $H(z,z')=\sum_i z_i\overline{z'_i}/\mathrm{area}(B.L)$ on either
coordinate space, and $E_H=\operatorname{Im} H$. Write $d=D.\mathrm{host}.d$,
$w=(1+i\sqrt d)/2$, $q(x)=\sqrt d\,\operatorname{Re}H(x,x)/2$.
The one- and four-block multipliers $\alpha,\beta$ are both
$\exp(\pi i q)$ on their respective spaces. We derive the ten fields
of \nolinkurl{StdData} in the order needed for the construction.
The ideal description of $P$ is \nolinkurl{FL.lam_mem}
(\path{GlobalTheory/GL3/PLinkThU.lean:528--540}); the inverse ideal is
\nolinkurl{AA0} (\path{LocalRingCardinality/DataHost.lean:153--154}).
See \nolinkurl{IsoKG.isom}, \nolinkurl{lamMM}, \nolinkurl{Base.gC_spec},
\nolinkurl{Base.StdOK} (\path{KernelLattices/LateralTheory.lean:19--33,71--78,119--121,128--155}),
\nolinkurl{alphaQ}, \nolinkurl{StdData}
(\path{QuadraticSemicharacters/ProductThetaMem.lean:16--20,52--64}).

\paragraph{The two period signs.}\phantomsection\label{r16:std-char}
The element $w$ is the image of the host quadratic integer
\nolinkurl{AR.om}, so both $P$ and $M$ are stable under $w$.
For $P$ this follows by multiplying an ideal element by an integer of
$K$. For $M$ apply the same identity to $p$ in each additive generator. The field
$B.t_0$ says that $\omega=B.\omega_0\in P$ and that
$w\omega/2,(1-w)\omega/2\notin P$. Stability gives
$\omega/2\notin P$, proving \nolinkurl{per} and \nolinkurl{half}.
The two nonzero classes of $\omega,w\omega$ in $P/2P$ are distinct
because their difference is nonzero there. They form a basis of the
two-dimensional $\mathbf F_2$-space. See \nolinkurl{SO.hst},
\nolinkurl{SO.half} (\path{GlobalTheory/GL3/PStdOK.lean:207--227}),
\nolinkurl{FourthRootStd.theta_eq}, \nolinkurl{FourthRootStd.lamMM_smul}
(\path{SatakeParameters/SpanRangeEmb4EqTop.lean:484--509}) and
\nolinkurl{SC.gen} (\path{GlobalTheory/GL3/PStdChar.lean:254--272}).

Here is the quadratic calculation, including the odd-integer parameter
of \nolinkurl{quadMultiplierSemicharacter}. On a Hermitian space linear
in the first argument let an additive lattice $J$ have integral $E_H$ and
be stable under $w_b=(b+i\sqrt d)/2$ for an odd integer $b$. For $x,y\in J$
write $e=E_H(x,y)$ and $k=E_H(w_bx,y)$, both integers. Taking the
imaginary part of $w_b H(x,y)$ gives
$\sqrt d\,\operatorname{Re}H(x,y)=2k-be$. At $x=y$ this gives
$q(x)=E_H(w_bx,x)\in\mathbf Z$. At $(x,y)$ the cross term of
$q(x+y)$ is $2k-be$, congruent to $e$ modulo 2. Consequently
\[
 \exp(\pi i q(x+y))=\exp(\pi i q(x))\exp(\pi i q(y))
                      \exp(\pi i E_H(x,y)).
\]
Each of the first two factors is a sign. This supplies the two
integrality statements and both clauses of the semicharacter law at
the general parameters of the Lean lemma
(\path{CayleyGaussianIntegrals/MatrixFormFourfoldBlockFirstBlockEmb.lean:178--232}).

For $J=P$ take $b=1$. In a period basis $E_H(x,y)$ is the coordinate
determinant up to orientation, since the area is the absolute
determinant of the basis. Thus $q(\omega)=E_H(w\omega,\omega)$ is an
odd integer by the basis just found modulo 2. Also
$q(w\omega)=((1+d)/4)q(\omega)$ is even since the host gives
$d\equiv7\pmod8$. The other planar multiplier is
$\epsilon(p)\exp(\pi i E_H(p,\omega))$, where $\epsilon(p)=1$
on $2P$ and $-1$ on its other cosets. To check its multiplication
law, use coordinates $(a,b)$ of the period basis. Then
$\epsilon(p)=(-1)^{a+b+ab}$, whose extra exponent on addition is the
determinant modulo 2. The exponential factor is an ordinary character.
At $\omega$ the two multipliers are both $-1$. At $w\omega$ they
are both $+1$, since $E_H(w\omega,\omega)$ is odd. Their quotient is a
sign character and therefore is trivial on $2P$. The basis test now
proves equality on $P$. Multiplying this equality over coordinates
proves \nolinkurl{betaStd} on $A$. This also proves \nolinkurl{SC.char_eq}
for an arbitrary period pair with the stated stability, torsion
exclusions and $d\equiv7\pmod8$. See \nolinkurl{thetaZeroCharacter}
(\path{EisensteinTranslationProducts/ThetaZeroTranslateProduct.lean:10--11}),
\nolinkurl{SC.q_odd}, \nolinkurl{SC.char_eq}, \nolinkurl{SC.stdChar_eq}
(\path{GlobalTheory/GL3/PStdChar.lean:115--155,198--249,274--359}).

\paragraph{Spanning and isometry.}\phantomsection\label{r16:std-isom}
Let $e(x)_{a,\tau}=\tau(x_a)$ for $x\in E^4$, the source's
\nolinkurl{emb4}. Choose an integral basis $b_j$ of $E$. The square
matrix $(\rho(b_j))_{\rho,j}$ of \emph{all} embeddings into $\mathbf C$
has nonzero determinant, since its squared determinant is the
discriminant of the basis of the separable field $E$. The $\tau$ are
distinct rows of this matrix. Therefore $(\tau(b_j))_{\tau,j}$ has full
row rank. Putting $b_j$ in a single one of the four slots of $x$
shows that $e(E^4)$ spans all four blocks over $\mathbf C$. For real spanning, let $c_j\in\mathbf R$ satisfy
$\sum_j c_j\tau(b_j)=0$ for every selected $\tau$. Conjugation gives
the same relation for each complementary embedding. Inverting the
full embeddings matrix now gives $c_j=0$ for all $j$.
There are $[E:\mathbf Q]=2|\mathrm{Tau}|$ integral basis vectors,
the real dimension of the selected embedding space. For the discriminant
identity of a basis of a finite separable extension see J.~S.~Milne,
\emph{Algebraic Number Theory}, Proposition~2.26
(\url{https://www.jmilne.org/math/CourseNotes/ANT301.pdf}).
See \nolinkurl{FL.linInd_embP}, \nolinkurl{FL.span_embP}
(\path{GlobalTheory/GL3/PLinkThU.lean:215--264}) and
\nolinkurl{FourthRootStd.span_range_emb4_eq_top}
(\path{SatakeParameters/SpanRangeEmb4EqTop.lean:256--295}).

The field \nolinkurl{Base.gC_spec} sends $e(x)$ to
$(\phi_0((\gamma_0x)_v))_v$. By the defining $K$-isometry of $\gamma_0$,
both Hermitian pairings of $e(x),e(y)$, before and after $g$, are
\[
 \frac{1}{\mathrm{area}(B.L)}\,
 \phi_0\!\left(\sum_a\operatorname{Tr}_{E/K}(x_a\overline{y_a})\right).
\]
Indeed the relative trace is the sum over the embeddings above
$\phi_0|_K$, and every embedding takes host conjugation to complex
conjugation. Since both sides are sesquilinear and $g$ is complex-linear,
the equality extends first in $x$, then in $y$, to all vectors.
The four-block matrix has diagonal blocks $H$ and zero off them, giving
\nolinkurl{isom}. Since $q(0)=0$ and $q$ is additive across orthogonal
blocks, $\alpha(0)=1$ and $\beta(gz)=\prod_a\alpha(z_a)$ for every
$z$, which give \nolinkurl{alpha0} and \nolinkurl{betaAlpha}.
See \nolinkurl{FL.embA_bij}, \nolinkurl{FL.sum_tau_trace}
(\path{GlobalTheory/GL3/PHfTrace.lean:26--70}),
\nolinkurl{Base.isomOK}
(\path{KernelLocalFactors/IsAdmissibleFlagPairOfUnitriangularLeftAutomorph.lean:748--798}),
\nolinkurl{FourthRootStd.isomOK} (\path{SatakeParameters/SpanRangeEmb4EqTop.lean:297--347})
and \nolinkurl{quadMultiplierRules}
(\path{CayleyGaussianIntegrals/MatrixFormFourfoldBlockFirstBlockEmb.lean:234--242}).

\paragraph{The image lattice and its dual.}\phantomsection\label{r16:std-dual}
For arbitrary ranks $(N,N')$ the standard lattices over
$\mathcal O_E$ have as their homomorphism lattice all integral
$N'\times N$ matrices. Images of the coordinate basis test integrality of each
entry, and the elementary matrices generate every such map. The
product of that module with $I$ consists of the matrices with entries
in $I$. One inclusion follows entry by entry, and $y$ times an
elementary matrix gives an arbitrary single entry $y\in I$ for the
converse. Thus a generator of the period lattice, with a single
coefficient $p\in P$, has each entry in $M$. Conversely every generator
$p(\tau y)_\tau$ of $M$ in a single matrix entry belongs to the period
lattice. Addition over all entries gives
$\mathrm{La}(S_N,S_{N'})=M^{N'\times N}$, including the empty matrix
case. This is \nolinkurl{Base.lat0OK} at arbitrary ranks. Also the
absolute determinant of any two periods is an integral multiple of the
basis area. The least positive such determinant is the basis area,
attained on the basis itself. This identifies the area of the additive
period lattice used in $\mathrm{La}$ with the area in $H$.
See \nolinkurl{FourthRootStd.Lat0}, \nolinkurl{FourthRootStd.areaOK}
(\path{SatakeParameters/SpanRangeEmb4EqTop.lean:103--129,169--249}),
\nolinkurl{Base.lat0OK} and \nolinkurl{Base.areaOK}
(\path{KernelLocalFactors/IsAdmissibleFlagPairOfUnitriangularLeftAutomorph.lean:318--319,395--439}).

The coordinate vectors generate $\mathcal O_E^2$ and span $E^2$ over $E$.
Pairing with them tests integrality, and conjugation preserves
$\mathcal O_E$, so this lattice is self-dual. Both ranks of the pair
$(2,2)$ are even. This verifies admission of the standard pair. The base period is nonzero, since
$\delta=0$ would give $P=0$, contrary to the period basis. Apply the
self-duality clause of the confirmed lattice law to this admitted pair.
Transporting it by the area and lattice identifications just proved
and reindexing the four entries gives
$(M^4)^\#=M^4$ for the alternating pairing of the block form.
Here $J^\#=\{z:E_H(z,y)\in\mathbf Z\ \text{for every }y\in J\}$.
If $u$ is an isometry, $z$ pairs integrally with every $uy$ exactly
when $u^{-1}z$ pairs integrally with every $y$. Applying this to
$g$ gives $(A')^\#=A'$, the field \nolinkurl{selfDual}.
It implies integrality of $E_H$ on $A'$. The lattice $M^4$ is
$w$-stable and $g$ is complex-linear, so $A'$ is also $w$-stable.
Applying the quadratic calculation at $b=1$, now on $A'$, proves
\nolinkurl{betaSemi}.
See \nolinkurl{Base.latAll} (Section~\ref{r14:lat-indices},
\path{KernelLocalFactors/IsAdmissibleFlagPairOfUnitriangularLeftAutomorph.lean:623--644}),
\nolinkurl{SO.selfDual}, \nolinkurl{SO.betaSemi}
(\path{GlobalTheory/GL3/PStdOK.lean:86--119,229--249}) and
\nolinkurl{FourthRootStd.selfDual}, \nolinkurl{FourthRootStd.betaSemi}
(\path{SatakeParameters/SpanRangeEmb4EqTop.lean:32--81,351--399,592--611}).
The latter implementation calls \nolinkurl{selfDual_La} directly. Its
forward dual test uses \nolinkurl{self_dual_lattice}, and integrality on
two lattice points follows by taking the imaginary part of
\nolinkurl{hermitian_form_integral_values}. These are the two ingredients
used for the same clause of \nolinkurl{Base.latAll}
(Sections~\ref{r14:lat-integral} and~\ref{r14:lat-dual}).

\paragraph{Two integral bounds and the indices.}\phantomsection\label{r16:std-indices}
Select $k>0$ with $kI\subset\mathcal O_E$. Indeed a nonzero integral
denominator for the fractional ideal $I$ divides the positive integer
given by the absolute norm of its principal ideal. For $a\in\mathrm{Fin}(4)$
and coordinate $v$ let $f_{av}(t)=(\gamma_0(e_a t))_v\in K$.
There are finitely many $f_{av}$ and finitely many integral basis
vectors of $\mathcal O_E$. Choose a single positive integer $m$ clearing
all their images into $\mathcal O_K$. By additivity it works on every
$t\in\mathcal O_E$. On $y\in I^4$ this gives
\[
 (mk)(\gamma_0 y)_v=\sum_a m f_{av}(ky_a)\in\mathcal O_K .
\]
The vectors $p\,e(y)$ with $p\in P$, $y\in I^4$ additively generate
$M^4$, as seen by putting a one-block generator in each separate
slot. Multiplying $g(p\,e(y))$ by $mk$ therefore gives a vector in
$A$, by \nolinkurl{Base.gC_spec} and the $\mathcal O_K$-stability of
$P$. This proves $mk A'\subset A$. This forward bound is also the
one used when clearing torsion shifts in the fourfold theta identity.
See \nolinkurl{FL.exists_int_clearing_order_denominators},
\nolinkurl{FL.one_mem_inverse_halfDifferent}
(\path{GlobalTheory/GL3/PLinkThU.lean:387--412}),
\nolinkurl{SO.denom}, \nolinkurl{SO.up}
(\path{GlobalTheory/GL3/PStdOK.lean:39--70,251--301}), and
\nolinkurl{FourthRootStd.up}
(\path{SatakeParameters/SpanRangeEmb4EqTop.lean:510--547,613--663}).

Conversely take a positive integer $h$ clearing the finitely many
entries of $\gamma_0^{-1}(e_v)$ into $\mathcal O_E$ for every $v$,
using the standard $K$-coordinate vectors $e_v$.
Since $\mathfrak b$ is a nonzero integral ideal,
$\mathcal O_E\subset\mathfrak b^{-1}=I$. Thus
$p\,e(h\gamma_0^{-1}(e_v))\in M^4$ for $p\in P$ and its $g$-image
is $hp e_v$. These generate $hA$, whence $hA\subset A'$.
In particular
\[
 hA\subset A\cap A'\subset A,A'\subset (mk)^{-1}A.
\]
A product of copies of the rank-two period lattice is finitely generated
over $\mathbf Z$. The quotient $(mk)^{-1}A/hA$ is finite, so
both $A/(A\cap A')$ and $A'/(A\cap A')$ are finite. This proves
nonvanishing of the two relative indices. With the source convention
\nolinkurl{relIdx}$(A,A')=[A':A\cap A']$ is
\nolinkurl{idx₁}. The reversed index is \nolinkurl{idx₂}.
See \nolinkurl{SO.down}, \nolinkurl{SO.idx₁}, \nolinkurl{SO.idx₂}
(\path{GlobalTheory/GL3/PStdOK.lean:303--336}),
\nolinkurl{FourthRootStd.down}, \nolinkurl{FourthRootStd.idx₁},
\nolinkurl{FourthRootStd.idx₂}
(\path{SatakeParameters/SpanRangeEmb4EqTop.lean:665--698}) and
\nolinkurl{relIdx} (\path{QuadraticSemicharacters/IsSemiChar.lean:63}).

All ten fields are now attached to the specified base. In the
\nolinkurl{SO.stdOK} implementation \nolinkurl{Base.stdOK_of'} inserts
\nolinkurl{Base.isomOK}. In \nolinkurl{FourthRootStd.stdOK} the
corresponding isometry is an explicit argument of
\nolinkurl{FourthRootStd.stdOK_of}. The preceding derivation uses
the period, the $K$-isometry and its complex extension carried by
the base, and the independently established lattice duality. It
supplies the standard-data hypothesis before a fourth root or a
translated jet is chosen. Review A confirmed both Lean entry points
for this argument
(\path{KernelLocalFactors/IsAdmissibleFlagPairOfUnitriangularLeftAutomorph.lean:327--352,800--814},
\path{SatakeParameters/SpanRangeEmb4EqTop.lean:137--163,700--701},
\path{GlobalTheory/GL3/PStdOK.lean:338--339}).

% blueprint/ch_P_r16.tex lines 261 through 450; r16:ktheta, r16:ktheta-line, r16:ktheta-transport, r16:ktheta-root
\subsection{The fourth-root theta vector}\label{r16:ktheta}
Let $L=\mathbb Z\omega_1+\mathbb Z\omega_2$ be the lattice of a
period pair. Let $I$ be a finite index set, $V=\mathbb C^I$,
$h$ the matrix of the one-block form, and $\Lambda\subset V$ the
additive subgroup in the statement. Write $A=L^{4\times I}$ and
$M=g(\Lambda^4)$ in the standard four-block space. The ten hypotheses
of \nolinkurl{StdData} are the following. The period $\omega_0$ belongs
to $L$ and $t=\omega_0/2$ does not. The complex linear equivalence $g$
carries the four-block form to the standard one. The one-block multiplier
satisfies $\alpha(0)=1$. The common function $\beta$ equals the standard
product character on $A$, is a semicharacter on $M$, and satisfies
$\beta(gx)=\prod_{a=0}^3\alpha(x_a)$ on the whole space. Finally $M$ is
self-dual and both relative indices of $A,M$ are nonzero. Fix
$\kappa\in\mathbb C^\times$. Write $F_0$ for the normalized standard
product and $I_{AM}$ for the source's finite intertwiner. We will obtain
an entire one-block theta vector $\theta$ satisfying
\[
 \prod_{a=0}^3\theta(z_a)=\kappa (I_{AM}F_0)(g z).
\]
Every solution in the indicated theta space is nonzero and spans it.
The solutions differ by multiplication by a fourth root of unity.
See \nolinkurl{StdData}, \nolinkurl{kernel4}, \nolinkurl{ThetaVector}
(\path{QuadraticSemicharacters/ProductThetaMem.lean:52--78}) and
\nolinkurl{Ker.kThetaVector}
(\path{KernelLocalFactors/IsAdmissibleFlagPairOfUnitriangularLeftAutomorph.lean:23--29}).

\paragraph{The elliptic line and its coordinate products.}
\phantomsection\label{r16:ktheta-line}
Put $A_L=|\Im(\bar\omega_1\omega_2)|$,
$H_L(z,w)=z\bar w/A_L$ and $E_L=\Im H_L$.
For $s(\omega)=1$ on $2L$ and $-1$ on $L\setminus 2L$,
the modified sigma function $f(z)=\sigma_L(z)\exp(-s_2 z^2/2)$
has simple zeros exactly at $L$ and satisfies
\[
 f(z+\omega)=s(\omega)
   \exp\!\left(\pi H_L(z,\omega)+\tfrac\pi2 H_L(\omega,\omega)\right)f(z).
\]
In Lean, \nolinkurl{wSigma} and \nolinkurl{wZeta} give DLMF's product
and sum away from their poles: with $z/0=0$ the zero lattice vector
contributes the factor $1$ to the product and the term $1/z$ to the sum.
The real-linear \nolinkurl{wEta} assigns $2\zeta(\omega_i/2)$ to each
full period by DLMF (23.2.11)--(23.2.12), and the quotient defining
\nolinkurl{s2} is the coefficient of $\omega$ obtained by solving
$\eta(\omega)=s_2\omega+c\bar\omega$ on the two generators
(\path{EisensteinTranslationProducts/SemiChar.lean:8--23}).
DLMF's full periods are twice its half-periods. Its sigma
multiplier for $m\omega_1+n\omega_2$ has sign $(-1)^{m+n+mn}=s(\omega)$.
The Legendre relation gives the real-linear decomposition of the full
zeta increment $\eta(\omega)=s_2\omega+\pi\bar\omega/A_L$
(use a positive ordering of the generators). Substitution in the
quasiperiod formula cancels the $s_2$ term. The exponential factor
preserves the zeros and their multiplicities. The only lattice
hypothesis here is real linear independence of the two periods.
See NIST Digital Library of Mathematical Functions, \S23.2(ii)
and (23.2.14), (23.2.17), \url{https://dlmf.nist.gov/23.2};
\nolinkurl{sgStar}, \nolinkurl{semiChar} and \nolinkurl{th0}
(\path{EisensteinTranslationProducts/SemiChar.lean:18--34}),
\nolinkurl{etaSplit}, \nolinkurl{sgStarQuasi}, \nolinkurl{sgStarBasic}
(\path{EisensteinTranslationProducts/QuasiPeriodConstUnique.lean:10--24}).

For a Hermitian form $H$ write
$U_xF(z)=F(z+x)\exp(-\pi H(z,x)-\pi H(x,x)/2)$.
The normalized translate $e_t=U_tf/(U_tf)(0)$ exists since $t\notin L$,
and $e_t(0)=1$. Its multiplier on $L$ is
$s(\omega)\exp(\pi i E_L(\omega,\omega_0))$, a sign since
$E_L(\omega,\omega_0)\in\mathbb Z$. This follows by commuting the
two translations of $f$: the commutator first gives the phase with a minus sign, which has the same value because its exponent is an integral multiple of $\pi i$. Its only zeros are simple at $L-t$.
If $F$ is entire with that same multiplier, $F/e_t$ descends
meromorphically to $\mathbb C/L$, with at most one simple pole.
The integral of $(F/e_t)\,dz$ on the boundary of a translated
fundamental parallelogram vanishes by cancellation of opposite sides.
The possible residue is zero, so the simple pole is removable. The
quotient is holomorphic on the compact torus and hence constant.
Evaluation at zero gives $F=F(0)e_t$.
This is the one-variable argument in \nolinkurl{thetaLine_of} and
\nolinkurl{even_theta_line}, with the half-period translation already
made in the denominator (\path{ThreeTermNumEqTwistedGammaOfSigmaZero.lean:380--422},
\path{AnalyticIdentityTheorems/EvenThetaSecond.lean:10--89};
\nolinkurl{thetaZeroCharacter}, \path{ThetaZeroTranslateProduct.lean:10--11}).

On a standard product with coordinate index set $J$ the candidate
$F_0(z)=\prod_{i\in J}e_t(z_i)$ is entire and equals one at zero.
The exponent for $U_x$ is the sum of the coordinate exponents, so the
product obeys the required multiplier. In any other vector $G$ in the
same space fix all coordinates but $i$. Periods supported in that
coordinate give the one-variable law. Consequently
$G(z)=G(z\text{ with }z_i=0)e_t(z_i)$ everywhere, including the
zeros of $e_t$. Iteration gives $G(z)=G(0)F_0(z)$, also for the empty
product. On the standard lattice the multiplier is a semicharacter:
the one-dimensional signs satisfy the addition law with factor
$\exp(\pi i E_L(x,y))$ and the factors multiply across coordinates.
The one-variable law follows by writing $x,y$ in the period basis:
the parity of the cross terms in $m+n+mn$ equals the parity of their
alternating determinant. The extra half-period phase is a character.
The standard lattice is self-dual since the alternating matrix in
the period basis has off-diagonal entry $\pm1$; testing on each
coordinate period detects both integral coordinates.
See \nolinkurl{StdSelfDual_proof}, \nolinkurl{StdSemiChar_proof},
\nolinkurl{product_theta_mem}, \nolinkurl{th_split}, \nolinkurl{ThStd_proof}
(\path{Th0CharMulOpUTh0HalfPeriodEq.lean:152--182,184--197,199--212,214--280,282--309}).

\paragraph{Transport of the nonzero line.}
\phantomsection\label{r16:ktheta-transport}
Let $K=A\cap M$. On either lattice write $D_x=\beta(x)U_x$.
Each lattice is self-dual, so $E(x,y)=\Im H(x,y)$ is integral for
$x,y$ on that lattice. Thus $\exp(-2\pi i E(x,y))=1$, and
the Weyl product and the semicharacter identity give $D_xD_y=D_{x+y}$
for $x,y$ in that lattice. Thus on an $A$-fixed function the finite sum
\[
 I_{AM}F=|M/K|^{-1/2}\sum_{y\in M/K}D_yF
\]
is independent of representatives, is $M$-fixed by reindexing, and is
entire if $F$ is entire. For $x\in A$ and $y\in M$ the commutator
is $\exp(-2\pi i\,\Im H(x,y))$. This gives a well-defined pairing
of $A/K$ and $M/K$ by the same integrality.
Its kernel in $M/K$ is zero by self-duality of $A$ and its kernel in
$A/K$ is zero by self-duality of $M$. Both quotients are finite by
the two index hypotheses. Each injects into the character group of
the other, so they have the same order $d$. In the unnormalized
double sum defining $I_{MA}I_{AM}F$, move each $D_x$ through $D_y$
until it fixes $F$. For fixed $y$ the resulting sum over $x$ is zero
off the identity coset: a translation by an $x_0$ on which its
character is nontrivial multiplies that sum by a scalar different from one. At the identity coset
the sum is $dF$. Thus the two factors $d^{-1/2}$ give
$I_{MA}I_{AM}F=F$. Interchanging the lattices gives the inverse on
$M$-fixed functions. In particular $I_{AM}F_0$ is nonzero; applying
$I_{MA}$ to an arbitrary vector and transporting its scalar multiple
of $F_0$ back shows that $I_{AM}F_0$ spans the moved space.
See \nolinkurl{heisenbergComm}, \nolinkurl{semicharacterHom}
(\path{CayleyGaussianIntegrals/MatrixFormFourfoldBlock.lean:36--46,125--139}),
\nolinkurl{InterReps_of}, \nolinkurl{InterFixed_of}, \nolinkurl{InterTh_of}
(\path{CayleyGaussianIntegrals/LocalComponentGL.lean:103--120,130--181}),
\nolinkurl{characterSum}, \nolinkurl{IndexEq_of}, \nolinkurl{InterInv_of}
and \nolinkurl{LineTransport_of}
(\path{CayleyGaussianIntegrals/LogarithmEquivalence.lean:155--199,201--254}).

Pullback by $g$ carries the moved space onto the theta space on
$\Lambda^4$: substituting $gz$ and $gx$ in $U$ preserves its
exponent by the isometry and replaces $\beta(gx)$ by the four
one-block factors. Composition with $g^{-1}$ gives the inverse,
and both maps are holomorphic. Therefore
$P(z)=\kappa(I_{AM}F_0)(gz)$ is a nonzero generator there.
See \nolinkurl{ThIsom_of} (\path{CayleyGaussianIntegrals/LocalComponentGL.lean:183--229})
and \nolinkurl{ThetaVector_of}
(\path{MatrixFormFourfoldBlockFirstBlockEmb.lean:250--302}).

\paragraph{Restriction and a single scalar fourth root.}
\phantomsection\label{r16:ktheta-root}
Choose $(u_0,u_1,u_2,u_3)$ with $P(u_0,u_1,u_2,u_3)\ne0$ and set
$q(z)=P(z,u_1,u_2,u_3)$. This is entire. Testing the four-block
translation law only at $(x,0,0,0)$ gives exactly the one-block law
for $q$, since the form is block-diagonal and $\alpha(0)=1$.
The tensor product of any four one-block vectors is a four-block
vector, by the same splitting of exponents. In particular
$q^{\otimes4}=cP$ for some $c\in\mathbb C$. Evaluation at
$(u_0,u_0,u_0,u_0)$ gives $c\ne0$ since $q(u_0)^4\ne0$.
Choose $r\in\mathbb C$ with $r^4=c^{-1}$. The entire function
$\theta=rq$ satisfies the identity on all four independent variables.

For any solution $v$, $P\ne0$ supplies a point $w$ with $v(w)\ne0$.
If $G$ is any one-block theta vector, $G\otimes v^{\otimes3}=bP$.
Evaluating at $(z,w,w,w)$ and canceling $v(w)^3$ gives $G=bv$
everywhere. In particular a second solution $v'$ is $bv$, and
$b^4P=P$ forces $b^4=1$. The argument only divides by nonzero
constants and remains valid on the zeros of $P$; it also applies when
$I$ is empty. See \nolinkurl{restrict_mem}, \nolinkurl{fourfoldTensorMem},
\nolinkurl{FourthRoot_of}, \nolinkurl{fourth_root}
(\path{MatrixFormFourfoldBlockFirstBlockEmb.lean:69--88,90--100,102--160}).

With $h=\mathrm{hOne}(L,I)$ and the two quadratic multipliers this
conclusion gives a witness for \nolinkurl{IsThM}. Consequently the
epsilon choice \nolinkurl{thM} has both defining properties and is
nonzero. At a base $B$ the specified $\kappa$ is nonzero because each
chosen root has a nonzero prescribed power and its integral powers
are nonzero. The source uses a finite-support product: for finite multiplicative support it is a product of nonzero factors, and for infinite support its defined value is one. Specialization supplies $B.\mathrm{thMM}$ and its
four-block identity. These deductions only use the preceding theta
laws and the fields of $\mathrm{StdOK}$.
See \nolinkurl{IsThM}, \nolinkurl{thM}, \nolinkurl{StdOK}
(\path{KernelLattices/LateralTheory.lean:20--33,149--157}),
\nolinkurl{thM_spec}, \nolinkurl{isThM_ne_zero}, \nolinkurl{Base.thMM_spec}
(\path{KernelLocalFactors/IsAdmissibleFlagPairOfUnitriangularLeftAutomorph.lean:75--85,182--193,216--222}).
The coefficient route uses a second implementation of this proposition,
\nolinkurl{GlobalTheory.Mem.Col2.ThetaVec.thetaVector}
(\path{IwasawaDecomposition/PeriodLatticeSesquilinearFormIsHerm.lean:437--498}).
Its proof uses its own standard-line, transport and fourth-root lemmas.
In the Pair cone it supplies \nolinkurl{WeilVac.thMM_spec}
(\path{IwasawaDecomposition/ExistsWeilOThZeroTriv.lean:33--36}).
The column-data copy \nolinkurl{Ker.Cd.kThetaVector} proves the same
statement outside the Pair cone
(\path{ColumnData/IsAdmissibleFlagPairOfUnitriangularLeftAutomorph.lean:28--30}).

% blueprint/ch_P_r16.tex lines 452 through 483; r16:pair-candidate
\subsection{Application to the translated kernel and the adjacent-site theorem}
\label{r16:pair-candidate}
Start with the actual base over the specified curve ideal and period
root. Section~\ref{r16:stdok} obtains its standard data from the
base and the reviewed lattice laws. Section~\ref{r16:ktheta} supplies
the fourth-root identity and entireness at these data and the base's
nonzero $\kappa$. The two formal specializations of the same epsilon
choice precede the translated-coefficient result:
\nolinkurl{Base.thMM_spec} feeds the entireness branch, and
\nolinkurl{WeilVac.thMM_spec} feeds the coefficient branch
(\path{KernelLocalFactors/IsAdmissibleFlagPairOfUnitriangularLeftAutomorph.lean:216--222},
\path{IwasawaDecomposition/ExistsWeilOThZeroTriv.lean:33--36}).
The coefficient argument in Section~\ref{r15:translated} first
translates the elliptic factor, then the four independent slots. It
applies the two finite intertwiners and the algebraic pair character
at the same admitted pair. This proves the coefficient clause of
\nolinkurl{ThLoc.algTAll} and, together with entireness, gives
\nolinkurl{TwV.translate_entire_algebraic}
(\path{GaloisGeometry/AlgTAll.lean:12--26},
\path{SatakeParameters/TranslateEntireAlgebraic.lean:15--18}).
In particular, the proof of either theta input never invokes the
conclusion about translated coefficients. The individual coefficients
are algebraic; the common number-field assertion for the full family
in Z4.2.13 remains a separate obligation.

These results supply the translated kernel field in the reviewed
\nolinkurl{pairOK_allSites} assembly, using the same adjacent tuples
as in Section~\ref{r15:pair}
(\path{SatakeParameters/KernelSiteLatticeTuplePairsAdmissible.lean:23--47}).
The AI review in round 16 requested edits E1--E7. After the changes, the AI review in round 17 accepted the written argument for Z4.2.04 at its stated scope.
The original input Z4.2 still has open components.

% Parent source SHA-256: 99a8461100a8be8176578b515f15d452958d95ac91f86554c8727ce0a98ee36e
% blueprint/ch_P_r14.tex lines 1132 through 1185; r14:kernel-lattice
\paragraph{The period forms on this pair.}\label{r14:kernel-lattice}
Put $A_C=\operatorname{area}(\delta\phi_0(\mathfrak c))$ and let
$H(z,w)=A_C^{-1}\sum_{\tau,i,j}z_{\tau ij}\overline{w_{\tau ij}}$,
$E_H(x,y)=\operatorname{Im} H(x,y)$ and
$q(x)=\frac{\sqrt d}{2}\operatorname{Re}H(x,x)$. Set
$M=\mathrm{aaOf}(D,\mathfrak c)\operatorname{Hom}_{\mathcal O_E}(L,L')$
inside the rational matrices and
$\Lambda=\delta\,\mathrm{emb}(M)$ with the same embeddings over $\phi_0|_K$.
Here products mean modules of finite sums. To compare $\Lambda$ with
$\mathrm{La}(D,B_0.\mathrm{tor},L,L')$, take $b_0\in\mathfrak c$,
$t_0\in\mathrm{AA0}$ and an $E$-linear map $X$ carrying $L$ into $L'$.
The image of $b_0t_0 X$ is
$\delta\phi_0(b_0)\,\mathrm{emb}(t_0\operatorname{mat}X)$.
This is a generator of the analytic period group. Multiplying an
$\mathcal O_E$-coefficient into $X$ preserves its lattice condition.
Conversely a matrix sending $L$ to $L'$ yields exactly such an $E$-linear
map by multiplication on column vectors. Expand the analytic period
and the integral matrix module in their generators and distribute.
Every analytic generator is a sum of these same elements. This proves
the equality of the two groups for this pair, including either rank
zero. See \nolinkurl{lamK_eq} in \path{InertCoprimeLemmas/RepComb.lean:188--254,265--333}.

The first clause of $B_0.\mathrm{LatAll}$ identifies $\Lambda$ with
its dual for $E_H$ by the lattice proof on p.~\pageref{r14:lat-indices},
confirmed in R14 review A. Hence $E_H(x,y)$ is an integer on $\Lambda^2$.
For the diagonal let $\omega=(1+i\sqrt d)/2$. The CM period ideal is
stable under $\omega$ since $\omega$ is the $\phi_0$-image of an
integer of $K$. Multiplying the period factor in each analytic generator
shows that $\omega x\in\Lambda$ for $x\in\Lambda$. Since $H(x,x)$ is real
and the first variable is complex linear, the same integrality gives
\[
             q(x)=E_H(\omega x,x)\in\mathbf Z .
\]
Consequently $\alpha(x)=\exp(\pi i q(x))$ on periods equals the analytic
sign $(-1)^{\operatorname{round}(q(x))}$. Applying the remaining
semicharacter clause of $\mathrm{LatOK}$ at $x,y,x+y$ transports its
multiplication law to $\alpha$ (p.~\pageref{r14:lat-indices}), giving both
parts of \nolinkurl{laFacts_of_adm}.

For \nolinkurl{lattice_integral_of_facts} in its stated generality,
take any of the fields, embeddings, scales and lattices in
\nolinkurl{LatIntOfFacts} (\path{KernelLattices/FirstKernelLawReading.lean:492--496}).
The form $H_V$ is Hermitian for a real scale, and $\Lambda=\Lambda^\vee$
gives $\Im H_V(x,y)\in\mathbf Z$ for $x,y\in\Lambda$.
The sign field of $\mathrm{LaFacts}$ asserts $\exp(\pi i q(x))=\pm1$:
squaring gives $\exp(2\pi i q(x))=1$, so its real argument $q(x)$
belongs to $\mathbf Z$. These three conclusions give $\mathrm{LatInt}$
for the general datum. The CM argument above establishes the sign
hypothesis at the particular admitted pair.
See \path{GL3/PStdOK.lean:121--135},
\path{SelfDualLatticeNorms/SelfDualComap.lean:174--182,211--216},
\path{GL3/LinkThetaStatementFamily.lean:82--105},
\path{InertCoprimeLemmas/RepComb.lean:335--362} and
\path{KernelLocalFactors/PeriodLatticeFinitelyGenerated.lean:1524--1529}.

% blueprint/ch_P_r14.tex lines 1187 through 1259; r14:kernel-finite
\paragraph{The three transversals and their weights.}\label{r14:kernel-finite}
Use $A,B$ for the embedded half-lattices. For each $Q$ put
$e_Q=1$ for $Q\in T$, $e_Q=-1$ for $Q\notin T$ with $\bar Q\in T$,
and $e_Q=0$ otherwise. Their rational modules are specified at $Q$
by $c^{-e_Q}M_Q$ and $c^{e_Q}M_Q$, respectively. Multiplication by
$c$ sends both into $M_Q$ and sends $M_Q$ into $c^{e_Q}M_Q$.
Intersecting the local conditions and then embedding gives
$cA,cB\subseteq\Lambda$ and $c\Lambda\subseteq B$.
The module $M$ is finitely generated: take common denominators for
a standard basis in a multiple of $L$, generators of $L'$ in a
multiple of the standard lattice, and generators of $\mathrm{aaOf}$.
These embed a nonzero scalar multiple of $M$ in an integral matrix
module over the Noetherian ring $\mathcal O_E$. Its additive image
$\Lambda$ is finitely generated since $\mathcal O_E$ is finite over
$\mathbf Z$. Thus $A$ and $B$ are subgroups of $c^{-1}\Lambda$ and
$\Lambda\cap A$ is a subgroup of $\Lambda$, all finitely generated.
The following quotients are killed by $c$, so are finite. Their orders
have no prime factors outside $S$, since any finite abelian group
generated by $h$ elements and killed by $c$ is a quotient of
$(\mathbf Z/c\mathbf Z)^h$.
\[
  \frac{A}{A\cap\Lambda},\qquad
  \frac{\Lambda\cap A}{B\cap\Lambda\cap A},\qquad
  \frac{B}{B\cap\Lambda}.
\]
Choose their transversals $S_X,S_A,S_B$ in the respective numerators
as selected by \nolinkurl{repsOf} in \nolinkurl{xiV}.
In particular $|S_A|$ is positive, including when the middle quotient
is trivial. A representative $t\in S_A$ already lies in $\Lambda$.
See \path{KernelLattices/XiIndCClass.lean:166--192,216--233} and
\path{KernelLocalFactors/PeriodLatticeFinitelyGenerated.lean:217--275,1660--1708,1714--1753}.

Extend $\Xi$ to the archimedean space by its value at the chosen
preimage under $\mathrm{embH}_\delta$ where one exists and by zero
elsewhere; denote this function by $\Phi$.
The definition of \nolinkurl{kerTh} is \nolinkurl{thC} with seed $e$;
\nolinkurl{thC} is \nolinkurl{xiV} evaluated at volume one.
With this choice the two volume factors in \nolinkurl{hatS} leave
exactly $1/|S_A|$. The character inverse in the hat is
$\exp(2\pi iE_H(t,y))$. In the outer translation the factor
from $U_xU_y$ is $\exp(-\pi iE_H(x,y))$: subtract the exponent in
$U_{x+y}$ from those of $U_xU_y$ to obtain
$-\frac\pi2(H(x,y)-H(y,x))$. Substitution in the two finite sums
therefore gives, pointwise,
\begin{align*}
 \mathrm{kerTh}_{L,L',\Xi}(z)
   &=\sum_{x\in S_X}\sum_{y\in S_B}r_{x,y}\,U_{x+y}e(z),\\
 r_{x,y}&=\frac{\alpha(x)\alpha(y)e^{-\pi iE_H(x,y)}}{|S_A|}
       \sum_{t\in S_A}\Phi(t-x)e^{2\pi iE_H(t,y)} .
\end{align*}
Only the middle transversal produces a denominator.
Indeed $cx,cy,ct\in\Lambda$, so $c^2q(x),c^2q(y)$ and
$c^2E_H(x,y),c^2E_H(t,y)$ are integers. The phases in the displayed
weight are roots of unity (their $2c^2$-th powers equal one). They have
norm one at every $p$-adic embedding, even when $p=2$. For $p\notin S$,
the positive integer $|S_A|$ is a $p$-adic unit.
All values $\Phi(t-x)$ are integral there by the site law or the
zero branch. The nonarchimedean triangle inequality gives
$|r_{x,y}|_p\leq1$ in every such embedding and the formula gives
$r_{x,y}\in\overline{\mathbf Q}$. Moreover $c(x+y)\in\Lambda$, so one
and the same positive $S$-supported integer works for all shifts.
More generally, in the abstract finite-sum assertion each transversal
member has a positive $S$-supported annihilator modulo the integral
lattice. Their finite product serves in place of $c$ in the phase and
shift calculations. There one takes $|S_A|$ positive and $S$-supported
and the integrality of the sampled values $\Phi(t-x)$ as hypotheses.
Indexing $S_X\times S_B$ by a finite ordinal gives exactly
\nolinkurl{IsKerG}. This proves the Fourier expansion for
\nolinkurl{twist_lattice_kernel} and its specialization
\nolinkurl{theta_coset_kernel}. See \path{OrthogonalChartLevels/KerTh.lean:26--31},
\path{KernelLattices/XiIndCClass.lean:18--27,210--239},
\path{QuadraticSemicharacters/IsSemiChar.lean:19--25},
and \path{KernelLocalFactors/PeriodLatticeFinitelyGenerated.lean:126--196,276--290}.

% blueprint/ch_P_r14.tex lines 1261 through 1297; r14:kernel-jets
\paragraph{Taylor coefficients of the actual kernel.}\label{r14:kernel-jets}
For these self-dual lattices the rational span of $\Lambda$ is
contained in the saturation of the standard pair's period lattice,
as shown on page~\pageref{r8:Z4.2.04}. Thus every $x+y$ in the expansion
lies in that saturation. The translated standard vacuum has algebraic
normalized Taylor coefficients there by the genus-one and fourfold
computations on that page. Each of the two changes of lattice in
$B_0.\mathrm{thMoved}$ is a finite sum of translates with algebraic
scalar (semicharacter sign times inverse square root of an index),
as on page~\pageref{r7:Z4.2.04}. All its representatives are in the
standard saturation. Composition with an additional shift in this
saturation contributes a root of unity: multiply the two arguments into
the integral standard periods and use bilinearity of $E_H$.
The zero-index convention gives the zero sum. Multiplication by the
algebraic character factor at this pair gives $e$ and preserves the
conclusion. The function $e$ itself is entire: it belongs to the theta
space on the admitted pair by \nolinkurl{vector_bound_link}
(\path{GlobalTheory/GL3/ProvisionExtensionLink.lean:38--40});
the first theta condition is global complex differentiability.
Every $U_{x+y}e$ is therefore entire, by translation and multiplication
by the exponential of a complex affine-linear function of $z$.
In finite-dimensional complex coordinates, Osgood's lemma gives
analyticity on neighborhoods and hence smoothness to every complex
order. This is \nolinkurl{AnForm.Osgood.osgood} applied by
\nolinkurl{analyticOnNhd_of_differentiable_vc}
(\path{AnalyticOnNhdOfDifferentiableVc.lean:21--32,34--41}).
Taking the normalized derivative of multi-index $\beta$ at zero termwise
now gives $\sum_{x,y}r_{x,y}[z^\beta]U_{x+y}e\in\overline{\mathbf Q}$.
Finite linearity also gives holomorphy of the sum. This is the pair of
fields used in \nolinkurl{kerTh_entire_algCoef} and then in
\nolinkurl{pairOK_allSites}. The arithmetic norm bound above concerns the
finite weights; a bound for a normalized jet needs its own local
hypotheses. See \path{VacuumEntireAndNormalizedJetsAlgebraic.lean:262--316},
\path{GaloisGeometry/AlgTAll.lean:12--26},
\path{TranslateEntireAlgebraic.lean:15--18},
\path{AnalyticOnNhdOfDifferentiableVc.lean:34--41} and
\path{KerThEntireAlgCoef.lean:22--40}.

% blueprint/ch_P_r14.tex lines 1299 through 1331; r14:site-pair
\paragraph{Reassembling the site predicate.}\label{r14:site-pair}
Fix $j,j+1$ and valid tuples at these steps in the same dictionary.
Their even ranks and self-dual lattices give admission. Insert the
finite weight function of Section~\ref{r14:kernel-finite} into the
entire and algebraic coefficient statement of Section~\ref{r14:kernel-jets}.
The moved datum is identified by \nolinkurl{movedPairDatum_congr};
the pair scalar belongs to the algebraic closure by
\nolinkurl{latticeTuplePairScalar_mem_algebraicClosure}, and
\nolinkurl{chi_detId_and_chiFac_alg} supplies the two assertions
on its character factor. For each place the two frame determinants
are units by \nolinkurl{isUnit_det_frameMatrix_ofCharts}. Obtain the
two local spans by case distinction on their indices using
\nolinkurl{frame_localLattice_span_kernelLattice_zero} and
\nolinkurl{frame_localLattice_span_kernelLattice_succ}. This gives all
fields of the conditional site predicate
\nolinkurl{GlobalTheory.pairOK_allSites}
(\path{SatakeParameters/KernelSiteLatticeTuplePairsAdmissible.lean:23--47}).
The earlier rounds gave the determinant-character and algebraic pair
scalar calculation (\path{ch_P_r8_frames.tex}, section
\ref{r8:Z4.2.04}); the class dictionary is the conclusion of Z4.2.03.
At the R14 review, the substantive site inputs still requiring a hand
proof were \nolinkurl{KerSelfDual.isSDLat_latK_of_stdL} and
\nolinkurl{KerSp.locSp_latK_isPt}
(\path{IsSDLatLatKOfStdL.lean:19--34};
\path{KernelLatticeFirstFactorLatticeTuplePoint.lean:21--30}).
Beyond Z4.2.01, .03 and .05 their union contributes 2,136 weighted
core lines, including 2,004 theorem lines. The two index-specific span
calls and the moved datum \nolinkurl{xiM_datLoc}
(\path{MovedPairDatumDependsOnlyOnLevel.lean:13--28}) also require their
own reindexing steps. The finite expansion here supplies the entire
kernel and its algebraic Taylor coefficients.
The next subsection applies the related finite expansion to a chosen
tied reader under its own step hypotheses.

